\documentclass[12pt]{exam}
\usepackage{amsthm}
\usepackage{libertine}
\usepackage[utf8]{inputenc}
\usepackage[margin=1in]{geometry}
\usepackage{amsmath,amssymb}
\usepackage{multicol}
\usepackage[shortlabels]{enumitem}
\usepackage{siunitx}
\usepackage{cancel}
\usepackage{graphicx}
\usepackage{pgfplots}
\usepackage{listings}
\usepackage{tikz}
\usepackage{booktabs}
\usepackage{indentfirst}

\pgfplotsset{width=10cm,compat=1.9}
\usepgfplotslibrary{external}
\tikzexternalize

\newcommand{\class}{Research in CompMath} 
\newcommand{\examnum}{Homework 0} 
\newcommand{\examdate}{8/30/2026}
\newcommand{\timelimit}{}

\begin{document}
\pagestyle{plain}
\thispagestyle{empty}

\noindent
\begin{tabular*}{\textwidth}{l @{\extracolsep{\fill}} r @{\extracolsep{6pt}} l}
\textbf{\class} & \textbf{Name:} & \textit{Mary Kim}\\ 
\textbf{\examdate} &&\\
\end{tabular*}\\
\rule[2ex]{\textwidth}{2pt}
% -------

\begin{enumerate} [listparindent=1.5em]
\item 

\begin {enumerate}
\item Factor $2^{15} - 1 = 32,767$ into a product of two smaller positive integers.

$2^{15}-1\ =\ \left(2^{5}-1\right)\left(1+2^{b}+2^{2b}+...+2^{\left(a-1\right)b}\right)$

set $b = 5, a = 3$

$=\ \left(2^{5}-1\right)\left(1+2^{5}+2^{2(5)}+...+2^{\left(3-1\right)5}\right)$

$=\left(2^{5}-1\right)\left(\sum_{n=0}^{2}2^{5n}\right)$

$=31 \times 1057$

\item Find an integer $x$ such that $1 < x < 2^{32,767} - 1$ and $2^{32,767} - 1$ is divisible by $x$.  

$2^{32767}-1\ =\ \left(2^{31}-1\right)\left(1+2^{31}+2^{62}+...+2^{\left(1056\right)\left(31\right)}\right)$

This makes $2^{32767}-1$ a factor of $2^{17}-1$, making x = $2^{31}-1 = 2,147,483,647$.
\end {enumerate}

\item Make some conjectures about the values of $n$ for which $n - 1$ is prime or the values of $n$ for which $3n - 2n$ is prime.

For all integers $n > 1$, $3^{n} - 1$ is even and therefore not prime.

I noticed a pattern that if $3^{n} - 2^{n}$ is prime, then $n$ must be prime.

\begin{table}[htbp]
    \centering
    \caption{Primality Table for $n$, $3^n - 1$, and $3^n - 2^n$}
    \label{tab:prime_data}
    \begin{tabular}{cccccc}
    \toprule
    $n$ & $n$ Prime? & $3^n - 1$ & $3^n - 1$ Prime? & $3^n - 2^n$ & $3^n - 2^n$ Prime? \\
    \midrule
    2  & yes & 8         & no & 5         & yes \\
    3  & yes & 26        & no & 19        & yes \\
    4  & no  & 80        & no & 65        & no  \\
    5  & yes & 242       & no & 211       & yes \\
    6  & no  & 728       & no & 665       & no  \\
    7  & yes & 2186      & no & 2059      & no  \\
    8  & no  & 6560      & no & 6305      & no  \\
    9  & no  & 19682     & no & 19171     & no  \\
    10 & no  & 59048     & no & 58625     & no  \\
    11 & yes & 177146    & no & 175099    & no  \\
    12 & no  & 531440    & no & 527345    & no  \\
    13 & yes & 1594322   & no & 1586131   & no  \\
    14 & no  & 4782968   & no & 4766585   & no  \\
    15 & no & 14348906  & no & 14316139  & no  \\
    16 & no & 43046720  & no & 42981185  & no  \\
    17 & yes  & 129140162 & no & 129009091 & yes \\
    \bottomrule
    \end{tabular}
\end{table}

\item The proof of Theorem 3 gives a method for finding a prime number different from any in a given list of prime numbers.

\begin{enumerate}
\item Use this method to find a prime different from 2, 3, 5, and 7.

According to Theorem 3, there are infinitely many prime numbers.

If $m = p_1p_2\cdots p_n + 1$, then $m$ is a prime number.

$m = 2 \cdot 3 \cdot 5 \cdot 7 + 1 = 211$

\item Use this method to find a prime different from 2, 5, and 11.


With the same formula, you get

$m = 2 \cdot 5 \cdot 11 + 1 = 111$.
\end{enumerate}

\item Find five consecutive integers that are not prime.


The five consecutive integers 24, 25, 26, 27, and 28 are not prime.

I figured this out by counting up; I don't necessarily have a proof for this. 

\item Use the table in Figure I.1 and the discussion on p. 5 to find two more
perfect numbers.
\begin{enumerate}
\item If $2^{n} - 1$ is prime, then $2^{n-1}(2^{n} - 1)$ is a perfect number.

For $n = 5$, $2^{5} - 1 = 31$ is prime, so $2^{4}(2^{5} - 1) = \bf{496}$ is a perfect number.

For $n = 7$, $2^{7} - 1 = 127$ is prime, so $2^{6}(2^{7} - 1) = \bf{8128}$ is a perfect number.

\end{enumerate}

\item The sequence 3, 5, 7 is a list of three prime numbers such that each
pair of adjacent numbers in the list differ by two. Are there any more
such “triplet primes”?

    No, because the prime numbers become more sparse as they get larger. For example, the next prime number after 7 is 11, which is a difference of 4 from 7. The next prime number would be 13, which is a difference of 2 from 11, but the next prime number after that is 17, which is a difference of 4 from 13. This pattern continues, and there are no more triplet primes.

\item A pair of distinct positive integers (m, n) is called amicable if the sum of all positive integers smaller than n that divide n is m, and the sum of all positive integers smaller than m that divide m is n. Show that (220, 284) is amicable.

The positive integers smaller than 284 that divide 284 are 1, 2, 4, 71, and 142. The sum of these numbers is: $220$.

Similarily, the positive integers smaller than 220 that divide 220 are 1, 2, 4, 5, 10, 11, 20, 22, 44, 55, and 110. The sum of these numbers is: $284$.
\end{enumerate}

\end{document}